\documentclass[11pt,a4paper,oneside]{book}

\usepackage[T1]{fontenc}
\usepackage[utf8]{inputenc}
\usepackage[french]{babel}
\usepackage{lmodern}
\usepackage{microtype}
\usepackage{geometry}
\usepackage{graphicx}
\usepackage{booktabs}
\usepackage{amsmath}
\usepackage{amssymb}
\usepackage{hyperref}

\geometry{margin=3cm}

\title{Titre de la these}
\author{Prenom Nom}
\date{\today}

\begin{document}

\frontmatter

\pagenumbering{gobble}

\begin{titlepage}
  \centering
  {\Large Universite / Ecole doctorale\par}
  \vspace{2cm}
  {\Huge\bfseries Titre de la these\par}
  \vspace{1cm}
  {\Large Manuscrit de these\par}
  \vspace{2cm}
  {\Large Prenom Nom\par}
  \vfill
  Directeur ou directrice de these : Prenom Nom\par
  Laboratoire : Nom du laboratoire\par
  \vspace{1cm}
  {\large \today\par}
\end{titlepage}

\cleardoublepage
\pagenumbering{roman}

\chapter*{Remerciements}

Texte des remerciements.

\chapter*{Resume}

Resume de la these en francais.

\chapter*{Abstract}

Thesis abstract in English.

\tableofcontents

\mainmatter

\chapter{Introduction}

Presenter le contexte, la problematique, les objectifs et le plan du manuscrit.

\chapter{Etat de l'art}

Synthese de la litterature et positionnement du sujet.

\chapter{Methodes}

Description des donnees, outils, hypotheses et protocoles experimentaux.

\section{Modele}

\begin{equation}
  y = f(x) + \varepsilon
\end{equation}

\chapter{Resultats}

Presentation des resultats principaux.

\begin{table}[h]
  \centering
  \caption{Exemple de tableau.}
  \begin{tabular}{lcc}
    \toprule
    Experience & Mesure 1 & Mesure 2 \\
    \midrule
    A & 1.2 & 3.4 \\
    B & 1.5 & 3.8 \\
    \bottomrule
  \end{tabular}
  \label{tab:exemple}
\end{table}

\chapter{Discussion}

Interpretation des resultats, limites et implications.

\chapter{Conclusion}

Bilan des contributions et perspectives.

\backmatter

\begin{thebibliography}{9}

\bibitem{exemple}
Auteur, A.
\newblock Titre de la reference.
\newblock \emph{Journal}, 1(1), 1--10, 2026.

\end{thebibliography}

\end{document}

